\begin{table*}[t]
\centering
\footnotesize
\setlength{\tabcolsep}{4pt}
\renewcommand{\arraystretch}{1.12}
\caption{Same meaning, different entity in \textcolor{encolor}{English} and \textcolor{encolor}{Hindi} idioms; entities are \underline{underlined}. Cluster sizes (number of idioms, en/hi) show asymmetric lexicalization. Hindi and Arabic idioms are romanised with a literal gloss. Every idiom shown was re-checked against the shared meaning and unsupported ones removed (\S\ref{sec:analysis-meaning}).}
\label{tab:meaning-cases-enhi}
\begin{tabular}{@{}p{0.20\textwidth}p{0.37\textwidth}p{0.37\textwidth}@{}}
\toprule
Shared meaning\newline (en/hi idioms) & \textcolor{encolor}{English imagery} & \textcolor{encolor}{Hindi imagery} \\
\midrule
near death (1/12)
& within an inch of one’s life --- almost an \underline{inch} from one’s \underline{life}
& \textit{qabr ka munh jhaank kar aaye hain} peeking into the \underline{grave}'s \underline{mouth}; \textit{koi dam ka mehmaan hai} a \underline{guest} for a few \underline{breaths}; \textit{paanv gor mein latkaaye baithe hain} feet dangling in the \underline{grave} \\
very clever or cunning (1/12)
& as clever as a cartload of monkeys --- a \underline{cartload} of \underline{monkeys}
& \textit{gonu ojha ki billi} \underline{Gonu Ojha}'s \underline{cat} \\
sudden unexpected occurrence (1/12)
& crop up --- \textit{(no imagery)}
& \textit{aja-krpaniya nyaya} sudden \underline{misfortune} befalls someone \\
forced by necessity (1/12)
& Needs must when the devil drives. --- \underline{the devil} drives, so \underline{needs} must
& \textit{jevad se nada ghisna padta hai} \underline{waist cord} must be rubbed by \underline{the rope} \\
successive misfortunes (1/12)
& a string of bad luck --- a \underline{string} of \underline{bad luck}
& \textit{dhakke par dhakka lagta hai} one \underline{blow} after another lands; \textit{pehle se nipte nahin, doosra sir par taiyar} before one \underline{trouble} ends, another arrives \\
support during hard times (1/12)
& bridge over troubled waters --- a \underline{bridge} over \underline{troubled waters}
& \textit{tumhari aankh meri taang, chalo ghoom aave} your \underline{eye} my \underline{leg}, let's go around; \textit{bhai vahi jo vipad sahay} a \underline{brother} is one who helps in \underline{distress} \\
\bottomrule
\end{tabular}
\end{table*}

\begin{table*}[t]
\centering
\footnotesize
\setlength{\tabcolsep}{4pt}
\renewcommand{\arraystretch}{1.12}
\caption{Same meaning, different entity in \textcolor{encolor}{English} and \textcolor{zhcolor}{Arabic} idioms; entities are \underline{underlined}. Cluster sizes (number of idioms, en/ar) show asymmetric lexicalization. Hindi and Arabic idioms are romanised with a literal gloss. Every idiom shown was re-checked against the shared meaning and unsupported ones removed (\S\ref{sec:analysis-meaning}).}
\label{tab:meaning-cases-enar}
\begin{tabular}{@{}p{0.20\textwidth}p{0.37\textwidth}p{0.37\textwidth}@{}}
\toprule
Shared meaning\newline (en/ar idioms) & \textcolor{encolor}{English imagery} & \textcolor{zhcolor}{Arabic imagery} \\
\midrule
to leave a place (1/12)
& up sticks --- \textit{(no imagery)}
& \textit{anshatu min dhi'b, wa min 'ayri al-falaah} more active than an \underline{wolf} and a \underline{wild ass} \\
falling into a worse situation (1/12)
& out of the frying pan, into the fire --- leaving the \underline{frying pan} and entering the \underline{fire}
& \textit{kal-mustaghithi mina r-ramda'i bin-nar} like one seeking help from the \underline{burning ground} with \underline{fire}; \textit{kharij min il-hariqah qabluh il-ghurab zaghutuh} out of the \underline{fire}, then met by the \underline{crow} \\
never happen (1/12)
& That’ll be the day! --- \textit{(no imagery)}
& \textit{hatta yarji'a al-sahmu 'ala fawqihi} until the \underline{arrow} returns to its \underline{bow}; \textit{la af'alu kadha ma balla al-bahru sufatan, wa-ma anna fi al-furati qatratan} until the \underline{sea} wets a \underline{wool tuft}; until the \underline{Euphrates} has a \underline{drop} \\
praise or acknowledge someone (1/12)
& give someone their flowers --- present \underline{flowers} to \underline{someone}
& \textit{alayhi mina allahi lisanun salihah} a \underline{good tongue} from \underline{God} \\
hope for good fortune (1/12)
& fingers crossed --- \underline{fingers} \underline{crossed}
& \textit{shhal ma tal el-lil kaytlaa en-nhar} \textit{(no imagery)}; \textit{al-ghaym yahum wa-r-rabb rahum} the \underline{clouds} circle and the \underline{lord} is merciful; \textit{sabbiq al-khayr talqa al-khayr} precede \underline{good} and find \underline{good} \\
extremely fast (1/12)
& lightning-quick --- like \underline{lightning}, very fast
& \textit{asra'u mina al-rihi, wa mina al-barqi, wa mina al-isharati, wa mina al-jawabi, wa mina al-bayni, wa mina al-lamhi, wa mina al-tarfi, wa min lamh al-basari, wa min tarfi al-'ayni, wa min raj'i al-sada} faster than \underline{wind}, \underline{lightning}, and a \underline{glance}; \textit{amkhat min al-sahm} swifter than an \underline{arrow} \\
\bottomrule
\end{tabular}
\end{table*}

\begin{table*}[t]
\centering
\footnotesize
\setlength{\tabcolsep}{4pt}
\renewcommand{\arraystretch}{1.12}
\caption{Same meaning, different entity in \textcolor{zhcolor}{Chinese} and \textcolor{encolor}{Hindi} idioms; entities are \underline{underlined}. Cluster sizes (number of idioms, zh/hi) show asymmetric lexicalization. Chinese idioms are given with a literal gloss. Hindi and Arabic idioms are romanised with a literal gloss. Every idiom shown was re-checked against the shared meaning and unsupported ones removed (\S\ref{sec:analysis-meaning}).}
\label{tab:meaning-cases-zhhi}
\begin{tabular}{@{}p{0.20\textwidth}p{0.37\textwidth}p{0.37\textwidth}@{}}
\toprule
Shared meaning\newline (zh/hi idioms) & \textcolor{zhcolor}{Chinese imagery} & \textcolor{encolor}{Hindi imagery} \\
\midrule
extremely old and weak (1/12)
& \zh{树倒根摧} \underline{tree} falls, \underline{roots} crushed
& \textit{na munh men daant na pet men aant} no \underline{teeth} in \underline{mouth}, no \underline{intestines} in \underline{belly} \\
seize opportunity promptly (1/12)
& \zh{见之不取，思之千里} seeing it but not taking it, then thinking of \underline{thousands of miles} later
& \textit{uthi penth aathven din} the \underline{market} rises only on the \underline{eighth day}; \textit{chheenke ko uthna aur billi ka lapakna} the \underline{hanging shelf} rises and the \underline{cat} pounces \\
near death / old age (1/12)
& \zh{风烛草露} \underline{wind}-blown \underline{candle} and \underline{grass dew}
& \textit{koi dam ka mehman hai} a \underline{guest} of a few \underline{breaths}; \textit{dhola bal maut ko nijnani} \underline{white hair} means \underline{death} is near; \textit{na munh mein daant na pet mein aant} no \underline{teeth} in the \underline{mouth}, no \underline{intestines} in the \underline{belly} \\
attempting beyond one's strength (1/12)
& \zh{力不从心} \underline{strength} does not follow \underline{heart}
& \textit{eent ki paant dam madar} \underline{brick row} and \underline{strength} for a task; \textit{chhati par bal nahin, bhalu se ladai} \underline{chest hair} gone, fighting a \underline{bear}; \textit{tant-si deh, paon na hath, larn chalo suran ke sath} \underline{thread-thin body}, no \underline{feet} or \underline{hands}, go fight \underline{heroes} \\
on the verge of death (1/12)
& \zh{行将就木} about to enter \underline{wood}
& \textit{koi dam ka mehmaan hai} a \underline{guest} of one breath; \textit{paanv gor mein latkaae baithe hain} sitting with \underline{feet} dangling in the \underline{grave} \\
everyone has flaws (1/12)
& \zh{金无足赤，人无完人} \underline{gold} is never \underline{pure red}; \underline{people} are not \underline{perfect}
& \textit{kisi ki naak tedhi, kisi ki aankh tedhi} someone's \underline{nose} is crooked, someone's \underline{eye} is crooked \\
\bottomrule
\end{tabular}
\end{table*}

\begin{table*}[t]
\centering
\footnotesize
\setlength{\tabcolsep}{4pt}
\renewcommand{\arraystretch}{1.12}
\caption{Same meaning, different entity in \textcolor{zhcolor}{Chinese} and \textcolor{zhcolor}{Arabic} idioms; entities are \underline{underlined}. Cluster sizes (number of idioms, zh/ar) show asymmetric lexicalization. Chinese idioms are given with a literal gloss. Hindi and Arabic idioms are romanised with a literal gloss. Every idiom shown was re-checked against the shared meaning and unsupported ones removed (\S\ref{sec:analysis-meaning}).}
\label{tab:meaning-cases-zhar}
\begin{tabular}{@{}p{0.20\textwidth}p{0.37\textwidth}p{0.37\textwidth}@{}}
\toprule
Shared meaning\newline (zh/ar idioms) & \textcolor{zhcolor}{Chinese imagery} & \textcolor{zhcolor}{Arabic imagery} \\
\midrule
seize the favorable moment (1/12)
& \zh{趁热打铁} while \underline{the iron} is \underline{hot}, strike it
& \textit{sir wa qamarun laka} walk while \underline{the moon} is yours; \textit{in kan liqal'ik rih unfudhuh} if \underline{your pulling} has \underline{wind}, shake it off \\
deceptive speech and appearance (1/12)
& \zh{巧言令色} skillful words and pleasing face
& \textit{tahta jildi al-da'n qalbu al-adh'ub} under the \underline{sheep's} skin is the \underline{wolf's} heart \\
truth becomes evident (1/12)
& \zh{水落石出} \underline{water} recedes, \underline{stones} appear
& \textit{qad bayyana as-subhu lidhi aynayn} the \underline{morning} makes it clear to \underline{two-eyed person}; \textit{illi tihbal billil twild binnahar} what is conceived in the \underline{night} is born in the \underline{day} \\
exaggerate a matter (1/12)
& \zh{张大其事} make \underline{the matter} big
& \textit{zayy al-safafir 'aqlah w-ghalabah} \underline{shouting} and \underline{claims} exceed \underline{capacity}; \textit{ya'mal al-habbah qubbah} make a \underline{grain} into a \underline{dome}; \textit{al-'aza hara (sakhana) wal-mayyit kalb} the \underline{funeral} is hot, but the \underline{dead man} is a dog \\
old age (1/12)
& \zh{雪鬓霜毛} \underline{snow}-white \underline{temples} and \underline{frost}-like \underline{hair}
& \textit{akala alayhi al-dahru wa shariba} \underline{time} has eaten and drunk on him; \textit{malaka la tanbahu ya kalba al-dawm qad kunta nabban fama laka al-yawm} why do you not bark, \underline{dog} of \underline{eternity}; you were once barking; \textit{faht al-bahr} \underline{sea} breath, \underline{old age} and passing time \\
terrified by slightest thing (1/12)
& \zh{惊弦之鸟} a \underline{bird} startled by a \underline{string}
& \textit{zayyi buhrjaan al-tarbi'ah; sha'rat riih tihizzuh} like a \underline{bird} in \underline{spring}; a \underline{hair} of \underline{wind} shakes it \\
\bottomrule
\end{tabular}
\end{table*}

\begin{table*}[t]
\centering
\footnotesize
\setlength{\tabcolsep}{4pt}
\renewcommand{\arraystretch}{1.12}
\caption{Same meaning, different entity in \textcolor{encolor}{Hindi} and \textcolor{zhcolor}{Arabic} idioms; entities are \underline{underlined}. Cluster sizes (number of idioms, hi/ar) show asymmetric lexicalization. Hindi and Arabic idioms are romanised with a literal gloss. Every idiom shown was re-checked against the shared meaning and unsupported ones removed (\S\ref{sec:analysis-meaning}).}
\label{tab:meaning-cases-hiar}
\begin{tabular}{@{}p{0.20\textwidth}p{0.37\textwidth}p{0.37\textwidth}@{}}
\toprule
Shared meaning\newline (hi/ar idioms) & \textcolor{encolor}{Hindi imagery} & \textcolor{zhcolor}{Arabic imagery} \\
\midrule
blaming others for own faults (1/12)
& \textit{apni kai doosre ke sir} one's own \underline{k} on another's \underline{head}
& \textit{ja'alat ma biha biya wa-intalakat talmiz} she made what is in \underline{her} be in \underline{me}, then went off blaming; \textit{kayfa tubsiru al-qadha fi 'ayni akhika wa tadharu al-jidh'a al-mu'tarida fi 'aynika} how do you see the \underline{speck} in your brother's \underline{eye} and leave the \underline{beam} in your own eye; \textit{ya'qidu fi mithli al-sawab wa fi 'aynayhi mithlu al-jarra} he ties up in the likeness of \underline{right}, while in his two \underline{eyes} is like a \underline{jar} \\
cooperation in work (1/12)
& \textit{donon haath milkar kaam karte hain} \underline{two hands} work together
& \textit{id wahda ma tsaffafsh} a single \underline{hand} cannot clap; \textit{id lohda mattsaffafsh} a single \underline{hand} cannot clap; \textit{id ala id tisaad} \underline{hand} on \underline{hand} helps \\
refuses advice and persists (1/12)
& \textit{sab ki dava hai aadat ki dava nahin} \textit{(no imagery)}
& \textit{ra'sun lishawrin ma yutaru nu'ratuhu} \underline{head} for \underline{consultation}; its \underline{call} is not driven away; \textit{la dhanba li qad qultu lilqawmi istaqu} \textit{(no imagery)} \\
one's loss becomes another's gain (1/12)
& \textit{anaaj bikhre murghon khush} \underline{grain} scattered, \underline{chickens} happy
& \textit{id faraghet fi ukhta-ha} a \underline{hand} emptied into its \underline{sister}; \textit{tukoon fi idak tiksim lighayrak} what is in your \underline{hand} is divided for \underline{another} \\
neglecting own duties for others (1/12)
& \textit{apni bhumi par ghas jame, doosre ko karen godai} \underline{own land} grows \underline{grass}, yet one does \underline{weeding} for others
& \textit{gina nsa'udhu fi dafni abuh fat lina al-fas wimishi} we came to help him bury his \underline{father}, then the \underline{axe} was left and gone \\
evil returns to doer (1/12)
& \textit{aag khaayega so angaar hagega} one who \underline{fire} eats will \underline{embers} defecate
& \textit{man yazra'i ash-shawk la yahsud bihi al-'inaba} one who \underline{thorns} sows will not reap \underline{grapes}; \textit{ala ahliha tajni baraqishu} \underline{Baraqish} brought harm upon \underline{its people} \\
\bottomrule
\end{tabular}
\end{table*}
